\begin{workedexample}[Worked Example: Connector frequency ceiling]
A coaxial connector supports a single TEM mode only up to the frequency
where higher-order modes appear, set by its air-gap dimensions. Standard SMA is
usable to about $18\ \text{GHz}$; the smaller 2.92~mm (``K'') connector reaches
$\sim40\ \text{GHz}$, and 1.85~mm (``V'') $\sim65\ \text{GHz}$. Smaller air
dimensions push the moding frequency higher --- at the cost of power handling
and ruggedness.
\end{workedexample}

\begin{keytakeaways}
\begin{itemize}[leftmargin=1.4em,itemsep=2pt]
\item Precision connectors trade size for bandwidth: SMA $\sim$18 GHz, 2.92 mm $\sim$40 GHz, 1.85 mm $\sim$65 GHz.
\item Impedance ($50$ or $75\ \Omega$) and dielectric interface must match the cable.
\item Torque to spec and limit mating cycles; a damaged interface spoils return loss.
\end{itemize}
\end{keytakeaways}

\begin{exercises}
\item Which connector suits a $26\ \text{GHz}$ link: SMA or 2.92~mm?
\item Are SMA and 2.92~mm mechanically mateable?
\item Why does over-torquing or excessive mating degrade RF performance?
\end{exercises}

\furtherreading{Pozar~\cite{pozar}; Steer~\cite{steer}.}
